
\textbf{One sequence length for most experiments.} Every single-query
timing in Sections~\ref{sec:results-hardware} and
\ref{sec:profiling} uses one 118-residue sequence. The
sequence-length sweep behind Section~\ref{sec:scaling}'s fit covers
100 to 1000 residues, but on a different, rotated-alphabet input
family (Section~\ref{sec:methodology}) that we do not compare
timings against the 118-residue results.

\textbf{One TPU generation.} All TPU measurements use v5e; we have no
data on v6e, v7 (Ironwood), or the announced TPU 8t/8i generation
\citep{googletpuv8ann2026}, and make no claim about whether any
finding here holds on a different generation.

\textbf{Untrained weights.} AlphaFold2 runs with randomly initialized
parameters throughout (Section~\ref{sec:methodology}). This is
deliberate for a timing study and does not affect the compiled
graph's shape, but it means every result here is a statement about
compute, not about prediction quality, and we did not measure whether
timing with trained weights is the same.

\textbf{Asymmetric isolation.} The TPU ran as a dedicated Kubernetes
Job. CPU and GPU ran on shared Colab runtimes whose underlying
machine the provider chooses and can change between sessions
(Section~\ref{sec:rh-variability}). The session-to-session
variability we measured and report is therefore a property of the
CPU/GPU setup specifically, not evidence about TPU stability one way
or the other.

\textbf{The headline speedups are computed from baselines that did not
reproduce.} The 27.8$\times$ TPU-over-GPU figure divides the August GPU
steady state by the TPU one, and the 451$\times$ TPU-over-CPU figure
divides the August CPU steady state by it; both numerators moved on
rerun, the GPU one by roughly a factor of two and the CPU one by
1.6--1.7$\times$, gaps we could not explain
(Section~\ref{sec:rh-variability}). We keep the August figures because
they come from the same measurement campaign as the TPU numbers and no
TPU rerun exists to pair with the September values: substituting them
would compare across sessions in exactly the way this section's own
finding argues against. Our records do not, however, establish that the
three August backends ran in a single session: the result files carry
no timestamp or session identifier, archived notebook output dates the
CPU and GPU runs to 2026-08-08, and the TPU run is dated only by its
commit (Section~\ref{sec:meth-provenance}).
The consequence is that these ratios compare one chip against one T4
and two vCPUs respectively, on one measurement campaign whose baselines
later moved, and should not be read as portable hardware constants.

\textbf{Most TPU measurements are single runs.} Two exceptions: a
three-repeat check on the single-query configuration (coefficient of
variation under 1.2\% on every timed region), and the auto-mesh
experiment, which ran twice. Neither extends to the multi-chip
\texttt{pmap} or scaling-grid experiments. Access to the
TPU cluster ended with the course that provided it, so none of this
can be regenerated or repeated.

\textbf{Provenance gaps predating this project's audit.} The original
August baselines were produced by scripts that were never committed
to version control in the form that ran (Section~\ref{sec:methodology}),
so those specific numbers cannot be reproduced bit-for-bit from the
released repository. The September reruns close this gap for CPU and
GPU only. Several supporting figures, including the raw XLA profiler
trace behind Section~\ref{sec:profiling}, the vendor hardware
specifications in Section~\ref{sec:background}, and one AlphaFold3
Stanford-CPU timing, come from project documentation, with no raw data
file behind them in the repository.

\textbf{The measured software stack is only partly pinned.} On the TPU
Jobs only \texttt{jax[tpu]} was pinned, to 0.10.2;
\texttt{dm-haiku}, \texttt{tensorflow-cpu}, \texttt{numpy},
\texttt{biopython}, \texttt{ml\_collections} and \texttt{absl-py} were
installed unpinned at container start-up, and the AlphaFold2 commit was
not recorded (Section~\ref{sec:meth-software}). Re-running the same
scripts today may therefore resolve a different dependency set from the
one that produced these numbers, and any effect of that difference
cannot be separated from the hardware and session effects discussed
above.

\textbf{Cost figures rest on a single measurement per backend and one
pricing snapshot.} Section~\ref{sec:cost} prices each backend
from one steady-state throughput value, not from the session-level
distributions of Section~\ref{sec:rh-variability}, and uses on-demand
list prices from a single date. Neither changes qualitatively with
small input variation, but neither is a distribution.

\textbf{The scaling-law fit is an approximation, not a tight one.}
Section~\ref{sec:scaling} reports the fitted curve's actual deviation
from the measured grid directly, up to 32\% at the extremes of the
tested length range; with a single measurement per grid cell we do not
compute a confidence interval. A global two-exponent power law is a
convenient summary of this grid, not a precise model of it.

\textbf{AlphaFold3 results are effectively single runs per backend.}
Each of the five diffusion samples in Section~\ref{sec:alphafold3}
comes from one process invocation, not five independent runs, so the
CPU-versus-GPU divergence there is one comparison, not a distribution
of comparisons, and the documented cause we report for it
(Appendix~\ref{app:af3-repro}) is plausible but unconfirmed.
